% =====================================================================
% Paper 15 — GAP-3 SO(10) anomaly-cancellation auxiliary note
%   (PHYSICS ERROR REPAIRED: RHN singlet hypercharge Y=1/2 -> Y=0)
% Status: [DRAFT] [NEEDS_UPDATE] [PHYSICS_ERROR_REPAIRED]
% AUDIT-FLAG: AUDIT-2026-05-02-Wave7-Aux-Epoch-Overclaim
%   (Math314-AddD extension covering Wave 1/4/5 papers Paper-09..16).
%   Hostile-referee audit (2026-05-02) identified THREE issues:
%   (1) PHYSICS ERROR: original Step 3 stated "RHN singlet 1(5) has
%       hypercharge y=1/2", but the SO(10) right-handed neutrino in
%       16 -> 1(5) under SU(5)xU(1)_chi has SM hypercharge Y=0
%       (it is a SM gauge singlet by definition; the (5) is the
%       U(1)_chi charge, NOT U(1)_Y hypercharge). REPAIRED in place.
%   (2) Anomaly taxonomy needs precision: SU(2)^3 perturbative is
%       group-theoretically trivial (relevant point is Witten global
%       anomaly); Tr(U(1)_Y) is the mixed gravitational anomaly
%       coefficient. Refined.
%   (3) "verified to six-digit precision" -> "symbolically exact zero,
%       numerically cross-checked" (anomaly cancellation is exact
%       algebraic, not numerical estimate; per Math319 V1 acceptance
%       gate). Refined.
% Compile: pdflatex Paper-15.tex (×2 for refs)
% Canonical archive: Math157, Math281, Math319 V1, Math314-AddD
% =====================================================================
% --- TECT canonical preamble (see _shared/tect-preamble.tex) ----------
% =====================================================================
% TECT canonical LaTeX preamble (revtex4-2 / single-column / theorem-safe)
% =====================================================================
% Maintainer: Jusang Lee (jtkor@outlook.com)
% Binding from: 2026-05-06
% Purpose: a single source-of-truth preamble for every TECT paper
% (Paper-00 .. Paper-10 + Auxiliary notes). Designed to (a) eliminate
% font-corruption on pdflatex (T1 fontenc + lmodern), (b) make
% theorem-heavy mathematical-physics content render cleanly in a
% single-column layout (PRD class option, NOT PRL two-column), (c)
% provide consistent theorem numbering across all TECT papers via
% amsthm with a shared counter set, (d) make hyperref + cleveref
% cross-references resolve cleanly with the standard two-pass
% pdflatex compile.
%
% Usage in each Paper-NN.tex:
%   % =====================================================================
% TECT canonical LaTeX preamble (revtex4-2 / single-column / theorem-safe)
% =====================================================================
% Maintainer: Jusang Lee (jtkor@outlook.com)
% Binding from: 2026-05-06
% Purpose: a single source-of-truth preamble for every TECT paper
% (Paper-00 .. Paper-10 + Auxiliary notes). Designed to (a) eliminate
% font-corruption on pdflatex (T1 fontenc + lmodern), (b) make
% theorem-heavy mathematical-physics content render cleanly in a
% single-column layout (PRD class option, NOT PRL two-column), (c)
% provide consistent theorem numbering across all TECT papers via
% amsthm with a shared counter set, (d) make hyperref + cleveref
% cross-references resolve cleanly with the standard two-pass
% pdflatex compile.
%
% Usage in each Paper-NN.tex:
%   % =====================================================================
% TECT canonical LaTeX preamble (revtex4-2 / single-column / theorem-safe)
% =====================================================================
% Maintainer: Jusang Lee (jtkor@outlook.com)
% Binding from: 2026-05-06
% Purpose: a single source-of-truth preamble for every TECT paper
% (Paper-00 .. Paper-10 + Auxiliary notes). Designed to (a) eliminate
% font-corruption on pdflatex (T1 fontenc + lmodern), (b) make
% theorem-heavy mathematical-physics content render cleanly in a
% single-column layout (PRD class option, NOT PRL two-column), (c)
% provide consistent theorem numbering across all TECT papers via
% amsthm with a shared counter set, (d) make hyperref + cleveref
% cross-references resolve cleanly with the standard two-pass
% pdflatex compile.
%
% Usage in each Paper-NN.tex:
%   \input{../_shared/tect-preamble.tex}
%   \begin{document}
%   ...
%   \end{document}
%
% Build (every paper): `pdflatex Paper-NN && pdflatex Paper-NN`
% (the second pass resolves \ref{} / \cref{} forward references).
% A bash/PowerShell wrapper that automates this is at
% Docs/papers/papers/_shared/build-paper.sh / build-paper.ps1.
% =====================================================================

% ---- Document class --------------------------------------------------
% PRD class option of revtex4-2 with [reprint,onecolumn] gives the same
% APS heading / by-line / abstract style as PRL but in a wider single
% column that handles long display equations and theorem environments
% without column-break artefacts. nofootinbib keeps citations out of
% footnotes. The 11pt option restores standard font metrics; with T1
% fontenc + lmodern this gives non-bitmap glyphs.
\documentclass[%
  aps,prd,reprint,onecolumn,nofootinbib,11pt,%
  amsmath,amssymb,floatfix,longbibliography%
]{revtex4-2}

% ---- Encoding & fonts ------------------------------------------------
\usepackage[T1]{fontenc}        % proper 8-bit font encoding (no font corruption)
\usepackage[utf8]{inputenc}     % allow UTF-8 source
\usepackage{lmodern}            % scalable Latin Modern (replaces bitmap CM)
\usepackage{textcomp}           % extra text symbols
\usepackage{microtype}          % subtle kerning improvements

% ---- UTF-8 → LaTeX character mappings --------------------------------
% Maps the Unicode characters that occur in TECT body text (legacy from
% the chat-content auto-archival rule, CLAUDE.md §4) to LaTeX-safe
% commands. Without these declarations, T1+inputenc rejects the chars
% with "Unicode character X not set up for use with LaTeX".
% Adding declarations centrally (here) is preferable to per-file edits.
\DeclareUnicodeCharacter{00A7}{\S}              % §  (section sign)
\DeclareUnicodeCharacter{00B2}{\ensuremath{^{2}}}        % ²
\DeclareUnicodeCharacter{00B3}{\ensuremath{^{3}}}        % ³
\DeclareUnicodeCharacter{00B5}{\ensuremath{\mu}}         % µ (micro)
\DeclareUnicodeCharacter{00B7}{\ensuremath{\cdot}}       % · (middle dot)
\DeclareUnicodeCharacter{00D7}{\ensuremath{\times}}      % ×
\DeclareUnicodeCharacter{00E4}{\"a}             % ä
\DeclareUnicodeCharacter{00E9}{\'e}             % é
\DeclareUnicodeCharacter{00F6}{\"o}             % ö
\DeclareUnicodeCharacter{00FC}{\"u}             % ü
\DeclareUnicodeCharacter{010C}{\v{C}}           % Č
\DeclareUnicodeCharacter{010D}{\v{c}}           % č
\DeclareUnicodeCharacter{0394}{\ensuremath{\Delta}}      % Δ
\DeclareUnicodeCharacter{03A3}{\ensuremath{\Sigma}}      % Σ
\DeclareUnicodeCharacter{03A9}{\ensuremath{\Omega}}      % Ω
\DeclareUnicodeCharacter{03B1}{\ensuremath{\alpha}}      % α
\DeclareUnicodeCharacter{03B2}{\ensuremath{\beta}}       % β
\DeclareUnicodeCharacter{03B3}{\ensuremath{\gamma}}      % γ
\DeclareUnicodeCharacter{03B4}{\ensuremath{\delta}}      % δ
\DeclareUnicodeCharacter{03B5}{\ensuremath{\varepsilon}} % ε
\DeclareUnicodeCharacter{03BA}{\ensuremath{\kappa}}      % κ
\DeclareUnicodeCharacter{03BB}{\ensuremath{\lambda}}     % λ
\DeclareUnicodeCharacter{03BC}{\ensuremath{\mu}}         % μ
\DeclareUnicodeCharacter{03BD}{\ensuremath{\nu}}         % ν
\DeclareUnicodeCharacter{03C0}{\ensuremath{\pi}}         % π
\DeclareUnicodeCharacter{03C1}{\ensuremath{\rho}}        % ρ
\DeclareUnicodeCharacter{03C3}{\ensuremath{\sigma}}      % σ
\DeclareUnicodeCharacter{03C4}{\ensuremath{\tau}}        % τ
\DeclareUnicodeCharacter{03C6}{\ensuremath{\varphi}}     % φ
\DeclareUnicodeCharacter{03C7}{\ensuremath{\chi}}        % χ
\DeclareUnicodeCharacter{03C8}{\ensuremath{\psi}}        % ψ
\DeclareUnicodeCharacter{03C9}{\ensuremath{\omega}}      % ω
\DeclareUnicodeCharacter{2013}{--}              % – (en dash)
\DeclareUnicodeCharacter{2014}{---}             % — (em dash)
\DeclareUnicodeCharacter{2018}{`}               % ‘
\DeclareUnicodeCharacter{2019}{'}               % ’
\DeclareUnicodeCharacter{201C}{``}              % “
\DeclareUnicodeCharacter{201D}{''}              % ”
\DeclareUnicodeCharacter{2070}{\ensuremath{^{0}}}        % ⁰
\DeclareUnicodeCharacter{2071}{\ensuremath{^{i}}}        % ⁱ
\DeclareUnicodeCharacter{2122}{\textsuperscript{TM}}     % ™
\DeclareUnicodeCharacter{2126}{\ensuremath{\Omega}}      % Ω (ohm sign)
\DeclareUnicodeCharacter{210F}{\ensuremath{\hbar}}       % ℏ
\DeclareUnicodeCharacter{2115}{\ensuremath{\mathbb{N}}}  % ℕ
\DeclareUnicodeCharacter{2102}{\ensuremath{\mathbb{C}}}  % ℂ
\DeclareUnicodeCharacter{211D}{\ensuremath{\mathbb{R}}}  % ℝ
\DeclareUnicodeCharacter{2124}{\ensuremath{\mathbb{Z}}}  % ℤ
\DeclareUnicodeCharacter{211A}{\ensuremath{\mathbb{Q}}}  % ℚ
\DeclareUnicodeCharacter{2192}{\ensuremath{\to}}         % →
\DeclareUnicodeCharacter{21D2}{\ensuremath{\Rightarrow}} % ⇒
\DeclareUnicodeCharacter{2200}{\ensuremath{\forall}}     % ∀
\DeclareUnicodeCharacter{2203}{\ensuremath{\exists}}     % ∃
\DeclareUnicodeCharacter{2208}{\ensuremath{\in}}         % ∈
\DeclareUnicodeCharacter{2209}{\ensuremath{\notin}}      % ∉
\DeclareUnicodeCharacter{2227}{\ensuremath{\wedge}}      % ∧
\DeclareUnicodeCharacter{2228}{\ensuremath{\vee}}        % ∨
\DeclareUnicodeCharacter{2229}{\ensuremath{\cap}}        % ∩
\DeclareUnicodeCharacter{222A}{\ensuremath{\cup}}        % ∪
\DeclareUnicodeCharacter{2245}{\ensuremath{\cong}}       % ≅
\DeclareUnicodeCharacter{2248}{\ensuremath{\approx}}     % ≈
\DeclareUnicodeCharacter{2260}{\ensuremath{\neq}}        % ≠
\DeclareUnicodeCharacter{2264}{\ensuremath{\leq}}        % ≤
\DeclareUnicodeCharacter{2265}{\ensuremath{\geq}}        % ≥
\DeclareUnicodeCharacter{22C5}{\ensuremath{\cdot}}       % ⋅
\DeclareUnicodeCharacter{2295}{\ensuremath{\oplus}}      % ⊕
\DeclareUnicodeCharacter{2297}{\ensuremath{\otimes}}     % ⊗

% ---- Mathematics & symbols ------------------------------------------
\usepackage{amsmath,amssymb,bm,mathrsfs,mathtools}
\usepackage{amsthm}             % theorem environments (load BEFORE hyperref)

% ---- Theorem environments (shared counter set, TECT canonical) -------
% All theorems / lemmas / propositions / corollaries / remarks share a
% single counter so cross-references read consistently across a paper.
% Definitions get a separate counter (definitions are numbered in their
% own sequence by editorial convention).
\newtheorem{theorem}{Theorem}
\newtheorem{lemma}[theorem]{Lemma}
\newtheorem{proposition}[theorem]{Proposition}
\newtheorem{corollary}[theorem]{Corollary}

\theoremstyle{remark}
\newtheorem{remark}[theorem]{Remark}
\newtheorem{example}[theorem]{Example}

\theoremstyle{definition}
\newtheorem{definition}[theorem]{Definition}
\newtheorem{conjecture}[theorem]{Conjecture}
\newtheorem{hypothesis}[theorem]{Hypothesis}

% ---- Tables / floats / graphics --------------------------------------
\usepackage{booktabs}           % \toprule / \midrule / \bottomrule
\usepackage{array}
\usepackage{graphicx}
\usepackage{enumitem}           % \begin{itemize}[leftmargin=...] options
\usepackage{hyphenat}           % \hyp{} for breakable hyphens in \texttt
\usepackage{seqsplit}           % \seqsplit{} for long unbreakable strings
\sloppy                         % allow stretchier inter-word spacing to
                                % reduce overfull \hbox warnings on long URLs
                                % and long \texttt tokens (paragraph-level).
\emergencystretch=3em           % last-resort hbox stretch budget

% ---- Cross-references (hyperref last among the standard packages,
%      cleveref AFTER hyperref) ----------------------------------------
\usepackage[%
  colorlinks=true,%
  linkcolor=blue,%
  citecolor=blue,%
  urlcolor=blue,%
  breaklinks=true,%
  bookmarks=true,%
  bookmarksnumbered=true%
]{hyperref}
\usepackage[capitalise,nameinlink]{cleveref}

% ---- TECT-canonical author block macros -----------------------------
% (defined here so every paper has the same author / affiliation style)
\newcommand{\tectauthor}{%
  \author{Jusang Lee}%
  \email{jtkor@outlook.com}%
  \affiliation{Independent researcher, Republic of Korea}%
}

% ---- TECT-canonical archive footer ----------------------------------
\newcommand{\tectarchivefooter}{%
  \par\medskip
  \noindent\small\textit{Canonical archive}:
  \url{https://github.com/TECT-OpenPhysics/TECT}.\par%
}

% ---- TECT TODO / AUDIT-FLAG / HONEST-SCOPE inline marks --------------
% Used in drafts to highlight pending audit items without disturbing
% the body text style. Suppress via \tectsuppressmarks (define empty).
\providecommand{\tectauditflag}[1]{\textcolor{red}{\textbf{[AUDIT-FLAG: #1]}}}
\providecommand{\tecttodo}[1]{\textcolor{orange}{\textbf{[TODO: #1]}}}
\providecommand{\tecthonestscope}[1]{\textit{Honest scope: #1.}}

% =====================================================================
% End of TECT canonical preamble
% =====================================================================

%   \begin{document}
%   ...
%   \end{document}
%
% Build (every paper): `pdflatex Paper-NN && pdflatex Paper-NN`
% (the second pass resolves \ref{} / \cref{} forward references).
% A bash/PowerShell wrapper that automates this is at
% Docs/papers/papers/_shared/build-paper.sh / build-paper.ps1.
% =====================================================================

% ---- Document class --------------------------------------------------
% PRD class option of revtex4-2 with [reprint,onecolumn] gives the same
% APS heading / by-line / abstract style as PRL but in a wider single
% column that handles long display equations and theorem environments
% without column-break artefacts. nofootinbib keeps citations out of
% footnotes. The 11pt option restores standard font metrics; with T1
% fontenc + lmodern this gives non-bitmap glyphs.
\documentclass[%
  aps,prd,reprint,onecolumn,nofootinbib,11pt,%
  amsmath,amssymb,floatfix,longbibliography%
]{revtex4-2}

% ---- Encoding & fonts ------------------------------------------------
\usepackage[T1]{fontenc}        % proper 8-bit font encoding (no font corruption)
\usepackage[utf8]{inputenc}     % allow UTF-8 source
\usepackage{lmodern}            % scalable Latin Modern (replaces bitmap CM)
\usepackage{textcomp}           % extra text symbols
\usepackage{microtype}          % subtle kerning improvements

% ---- UTF-8 → LaTeX character mappings --------------------------------
% Maps the Unicode characters that occur in TECT body text (legacy from
% the chat-content auto-archival rule, CLAUDE.md §4) to LaTeX-safe
% commands. Without these declarations, T1+inputenc rejects the chars
% with "Unicode character X not set up for use with LaTeX".
% Adding declarations centrally (here) is preferable to per-file edits.
\DeclareUnicodeCharacter{00A7}{\S}              % §  (section sign)
\DeclareUnicodeCharacter{00B2}{\ensuremath{^{2}}}        % ²
\DeclareUnicodeCharacter{00B3}{\ensuremath{^{3}}}        % ³
\DeclareUnicodeCharacter{00B5}{\ensuremath{\mu}}         % µ (micro)
\DeclareUnicodeCharacter{00B7}{\ensuremath{\cdot}}       % · (middle dot)
\DeclareUnicodeCharacter{00D7}{\ensuremath{\times}}      % ×
\DeclareUnicodeCharacter{00E4}{\"a}             % ä
\DeclareUnicodeCharacter{00E9}{\'e}             % é
\DeclareUnicodeCharacter{00F6}{\"o}             % ö
\DeclareUnicodeCharacter{00FC}{\"u}             % ü
\DeclareUnicodeCharacter{010C}{\v{C}}           % Č
\DeclareUnicodeCharacter{010D}{\v{c}}           % č
\DeclareUnicodeCharacter{0394}{\ensuremath{\Delta}}      % Δ
\DeclareUnicodeCharacter{03A3}{\ensuremath{\Sigma}}      % Σ
\DeclareUnicodeCharacter{03A9}{\ensuremath{\Omega}}      % Ω
\DeclareUnicodeCharacter{03B1}{\ensuremath{\alpha}}      % α
\DeclareUnicodeCharacter{03B2}{\ensuremath{\beta}}       % β
\DeclareUnicodeCharacter{03B3}{\ensuremath{\gamma}}      % γ
\DeclareUnicodeCharacter{03B4}{\ensuremath{\delta}}      % δ
\DeclareUnicodeCharacter{03B5}{\ensuremath{\varepsilon}} % ε
\DeclareUnicodeCharacter{03BA}{\ensuremath{\kappa}}      % κ
\DeclareUnicodeCharacter{03BB}{\ensuremath{\lambda}}     % λ
\DeclareUnicodeCharacter{03BC}{\ensuremath{\mu}}         % μ
\DeclareUnicodeCharacter{03BD}{\ensuremath{\nu}}         % ν
\DeclareUnicodeCharacter{03C0}{\ensuremath{\pi}}         % π
\DeclareUnicodeCharacter{03C1}{\ensuremath{\rho}}        % ρ
\DeclareUnicodeCharacter{03C3}{\ensuremath{\sigma}}      % σ
\DeclareUnicodeCharacter{03C4}{\ensuremath{\tau}}        % τ
\DeclareUnicodeCharacter{03C6}{\ensuremath{\varphi}}     % φ
\DeclareUnicodeCharacter{03C7}{\ensuremath{\chi}}        % χ
\DeclareUnicodeCharacter{03C8}{\ensuremath{\psi}}        % ψ
\DeclareUnicodeCharacter{03C9}{\ensuremath{\omega}}      % ω
\DeclareUnicodeCharacter{2013}{--}              % – (en dash)
\DeclareUnicodeCharacter{2014}{---}             % — (em dash)
\DeclareUnicodeCharacter{2018}{`}               % ‘
\DeclareUnicodeCharacter{2019}{'}               % ’
\DeclareUnicodeCharacter{201C}{``}              % “
\DeclareUnicodeCharacter{201D}{''}              % ”
\DeclareUnicodeCharacter{2070}{\ensuremath{^{0}}}        % ⁰
\DeclareUnicodeCharacter{2071}{\ensuremath{^{i}}}        % ⁱ
\DeclareUnicodeCharacter{2122}{\textsuperscript{TM}}     % ™
\DeclareUnicodeCharacter{2126}{\ensuremath{\Omega}}      % Ω (ohm sign)
\DeclareUnicodeCharacter{210F}{\ensuremath{\hbar}}       % ℏ
\DeclareUnicodeCharacter{2115}{\ensuremath{\mathbb{N}}}  % ℕ
\DeclareUnicodeCharacter{2102}{\ensuremath{\mathbb{C}}}  % ℂ
\DeclareUnicodeCharacter{211D}{\ensuremath{\mathbb{R}}}  % ℝ
\DeclareUnicodeCharacter{2124}{\ensuremath{\mathbb{Z}}}  % ℤ
\DeclareUnicodeCharacter{211A}{\ensuremath{\mathbb{Q}}}  % ℚ
\DeclareUnicodeCharacter{2192}{\ensuremath{\to}}         % →
\DeclareUnicodeCharacter{21D2}{\ensuremath{\Rightarrow}} % ⇒
\DeclareUnicodeCharacter{2200}{\ensuremath{\forall}}     % ∀
\DeclareUnicodeCharacter{2203}{\ensuremath{\exists}}     % ∃
\DeclareUnicodeCharacter{2208}{\ensuremath{\in}}         % ∈
\DeclareUnicodeCharacter{2209}{\ensuremath{\notin}}      % ∉
\DeclareUnicodeCharacter{2227}{\ensuremath{\wedge}}      % ∧
\DeclareUnicodeCharacter{2228}{\ensuremath{\vee}}        % ∨
\DeclareUnicodeCharacter{2229}{\ensuremath{\cap}}        % ∩
\DeclareUnicodeCharacter{222A}{\ensuremath{\cup}}        % ∪
\DeclareUnicodeCharacter{2245}{\ensuremath{\cong}}       % ≅
\DeclareUnicodeCharacter{2248}{\ensuremath{\approx}}     % ≈
\DeclareUnicodeCharacter{2260}{\ensuremath{\neq}}        % ≠
\DeclareUnicodeCharacter{2264}{\ensuremath{\leq}}        % ≤
\DeclareUnicodeCharacter{2265}{\ensuremath{\geq}}        % ≥
\DeclareUnicodeCharacter{22C5}{\ensuremath{\cdot}}       % ⋅
\DeclareUnicodeCharacter{2295}{\ensuremath{\oplus}}      % ⊕
\DeclareUnicodeCharacter{2297}{\ensuremath{\otimes}}     % ⊗

% ---- Mathematics & symbols ------------------------------------------
\usepackage{amsmath,amssymb,bm,mathrsfs,mathtools}
\usepackage{amsthm}             % theorem environments (load BEFORE hyperref)

% ---- Theorem environments (shared counter set, TECT canonical) -------
% All theorems / lemmas / propositions / corollaries / remarks share a
% single counter so cross-references read consistently across a paper.
% Definitions get a separate counter (definitions are numbered in their
% own sequence by editorial convention).
\newtheorem{theorem}{Theorem}
\newtheorem{lemma}[theorem]{Lemma}
\newtheorem{proposition}[theorem]{Proposition}
\newtheorem{corollary}[theorem]{Corollary}

\theoremstyle{remark}
\newtheorem{remark}[theorem]{Remark}
\newtheorem{example}[theorem]{Example}

\theoremstyle{definition}
\newtheorem{definition}[theorem]{Definition}
\newtheorem{conjecture}[theorem]{Conjecture}
\newtheorem{hypothesis}[theorem]{Hypothesis}

% ---- Tables / floats / graphics --------------------------------------
\usepackage{booktabs}           % \toprule / \midrule / \bottomrule
\usepackage{array}
\usepackage{graphicx}
\usepackage{enumitem}           % \begin{itemize}[leftmargin=...] options
\usepackage{hyphenat}           % \hyp{} for breakable hyphens in \texttt
\usepackage{seqsplit}           % \seqsplit{} for long unbreakable strings
\sloppy                         % allow stretchier inter-word spacing to
                                % reduce overfull \hbox warnings on long URLs
                                % and long \texttt tokens (paragraph-level).
\emergencystretch=3em           % last-resort hbox stretch budget

% ---- Cross-references (hyperref last among the standard packages,
%      cleveref AFTER hyperref) ----------------------------------------
\usepackage[%
  colorlinks=true,%
  linkcolor=blue,%
  citecolor=blue,%
  urlcolor=blue,%
  breaklinks=true,%
  bookmarks=true,%
  bookmarksnumbered=true%
]{hyperref}
\usepackage[capitalise,nameinlink]{cleveref}

% ---- TECT-canonical author block macros -----------------------------
% (defined here so every paper has the same author / affiliation style)
\newcommand{\tectauthor}{%
  \author{Jusang Lee}%
  \email{jtkor@outlook.com}%
  \affiliation{Independent researcher, Republic of Korea}%
}

% ---- TECT-canonical archive footer ----------------------------------
\newcommand{\tectarchivefooter}{%
  \par\medskip
  \noindent\small\textit{Canonical archive}:
  \url{https://github.com/TECT-OpenPhysics/TECT}.\par%
}

% ---- TECT TODO / AUDIT-FLAG / HONEST-SCOPE inline marks --------------
% Used in drafts to highlight pending audit items without disturbing
% the body text style. Suppress via \tectsuppressmarks (define empty).
\providecommand{\tectauditflag}[1]{\textcolor{red}{\textbf{[AUDIT-FLAG: #1]}}}
\providecommand{\tecttodo}[1]{\textcolor{orange}{\textbf{[TODO: #1]}}}
\providecommand{\tecthonestscope}[1]{\textit{Honest scope: #1.}}

% =====================================================================
% End of TECT canonical preamble
% =====================================================================

%   \begin{document}
%   ...
%   \end{document}
%
% Build (every paper): `pdflatex Paper-NN && pdflatex Paper-NN`
% (the second pass resolves \ref{} / \cref{} forward references).
% A bash/PowerShell wrapper that automates this is at
% Docs/papers/papers/_shared/build-paper.sh / build-paper.ps1.
% =====================================================================

% ---- Document class --------------------------------------------------
% PRD class option of revtex4-2 with [reprint,onecolumn] gives the same
% APS heading / by-line / abstract style as PRL but in a wider single
% column that handles long display equations and theorem environments
% without column-break artefacts. nofootinbib keeps citations out of
% footnotes. The 11pt option restores standard font metrics; with T1
% fontenc + lmodern this gives non-bitmap glyphs.
\documentclass[%
  aps,prd,reprint,onecolumn,nofootinbib,11pt,%
  amsmath,amssymb,floatfix,longbibliography%
]{revtex4-2}

% ---- Encoding & fonts ------------------------------------------------
\usepackage[T1]{fontenc}        % proper 8-bit font encoding (no font corruption)
\usepackage[utf8]{inputenc}     % allow UTF-8 source
\usepackage{lmodern}            % scalable Latin Modern (replaces bitmap CM)
\usepackage{textcomp}           % extra text symbols
\usepackage{microtype}          % subtle kerning improvements

% ---- UTF-8 → LaTeX character mappings --------------------------------
% Maps the Unicode characters that occur in TECT body text (legacy from
% the chat-content auto-archival rule, CLAUDE.md §4) to LaTeX-safe
% commands. Without these declarations, T1+inputenc rejects the chars
% with "Unicode character X not set up for use with LaTeX".
% Adding declarations centrally (here) is preferable to per-file edits.
\DeclareUnicodeCharacter{00A7}{\S}              % §  (section sign)
\DeclareUnicodeCharacter{00B2}{\ensuremath{^{2}}}        % ²
\DeclareUnicodeCharacter{00B3}{\ensuremath{^{3}}}        % ³
\DeclareUnicodeCharacter{00B5}{\ensuremath{\mu}}         % µ (micro)
\DeclareUnicodeCharacter{00B7}{\ensuremath{\cdot}}       % · (middle dot)
\DeclareUnicodeCharacter{00D7}{\ensuremath{\times}}      % ×
\DeclareUnicodeCharacter{00E4}{\"a}             % ä
\DeclareUnicodeCharacter{00E9}{\'e}             % é
\DeclareUnicodeCharacter{00F6}{\"o}             % ö
\DeclareUnicodeCharacter{00FC}{\"u}             % ü
\DeclareUnicodeCharacter{010C}{\v{C}}           % Č
\DeclareUnicodeCharacter{010D}{\v{c}}           % č
\DeclareUnicodeCharacter{0394}{\ensuremath{\Delta}}      % Δ
\DeclareUnicodeCharacter{03A3}{\ensuremath{\Sigma}}      % Σ
\DeclareUnicodeCharacter{03A9}{\ensuremath{\Omega}}      % Ω
\DeclareUnicodeCharacter{03B1}{\ensuremath{\alpha}}      % α
\DeclareUnicodeCharacter{03B2}{\ensuremath{\beta}}       % β
\DeclareUnicodeCharacter{03B3}{\ensuremath{\gamma}}      % γ
\DeclareUnicodeCharacter{03B4}{\ensuremath{\delta}}      % δ
\DeclareUnicodeCharacter{03B5}{\ensuremath{\varepsilon}} % ε
\DeclareUnicodeCharacter{03BA}{\ensuremath{\kappa}}      % κ
\DeclareUnicodeCharacter{03BB}{\ensuremath{\lambda}}     % λ
\DeclareUnicodeCharacter{03BC}{\ensuremath{\mu}}         % μ
\DeclareUnicodeCharacter{03BD}{\ensuremath{\nu}}         % ν
\DeclareUnicodeCharacter{03C0}{\ensuremath{\pi}}         % π
\DeclareUnicodeCharacter{03C1}{\ensuremath{\rho}}        % ρ
\DeclareUnicodeCharacter{03C3}{\ensuremath{\sigma}}      % σ
\DeclareUnicodeCharacter{03C4}{\ensuremath{\tau}}        % τ
\DeclareUnicodeCharacter{03C6}{\ensuremath{\varphi}}     % φ
\DeclareUnicodeCharacter{03C7}{\ensuremath{\chi}}        % χ
\DeclareUnicodeCharacter{03C8}{\ensuremath{\psi}}        % ψ
\DeclareUnicodeCharacter{03C9}{\ensuremath{\omega}}      % ω
\DeclareUnicodeCharacter{2013}{--}              % – (en dash)
\DeclareUnicodeCharacter{2014}{---}             % — (em dash)
\DeclareUnicodeCharacter{2018}{`}               % ‘
\DeclareUnicodeCharacter{2019}{'}               % ’
\DeclareUnicodeCharacter{201C}{``}              % “
\DeclareUnicodeCharacter{201D}{''}              % ”
\DeclareUnicodeCharacter{2070}{\ensuremath{^{0}}}        % ⁰
\DeclareUnicodeCharacter{2071}{\ensuremath{^{i}}}        % ⁱ
\DeclareUnicodeCharacter{2122}{\textsuperscript{TM}}     % ™
\DeclareUnicodeCharacter{2126}{\ensuremath{\Omega}}      % Ω (ohm sign)
\DeclareUnicodeCharacter{210F}{\ensuremath{\hbar}}       % ℏ
\DeclareUnicodeCharacter{2115}{\ensuremath{\mathbb{N}}}  % ℕ
\DeclareUnicodeCharacter{2102}{\ensuremath{\mathbb{C}}}  % ℂ
\DeclareUnicodeCharacter{211D}{\ensuremath{\mathbb{R}}}  % ℝ
\DeclareUnicodeCharacter{2124}{\ensuremath{\mathbb{Z}}}  % ℤ
\DeclareUnicodeCharacter{211A}{\ensuremath{\mathbb{Q}}}  % ℚ
\DeclareUnicodeCharacter{2192}{\ensuremath{\to}}         % →
\DeclareUnicodeCharacter{21D2}{\ensuremath{\Rightarrow}} % ⇒
\DeclareUnicodeCharacter{2200}{\ensuremath{\forall}}     % ∀
\DeclareUnicodeCharacter{2203}{\ensuremath{\exists}}     % ∃
\DeclareUnicodeCharacter{2208}{\ensuremath{\in}}         % ∈
\DeclareUnicodeCharacter{2209}{\ensuremath{\notin}}      % ∉
\DeclareUnicodeCharacter{2227}{\ensuremath{\wedge}}      % ∧
\DeclareUnicodeCharacter{2228}{\ensuremath{\vee}}        % ∨
\DeclareUnicodeCharacter{2229}{\ensuremath{\cap}}        % ∩
\DeclareUnicodeCharacter{222A}{\ensuremath{\cup}}        % ∪
\DeclareUnicodeCharacter{2245}{\ensuremath{\cong}}       % ≅
\DeclareUnicodeCharacter{2248}{\ensuremath{\approx}}     % ≈
\DeclareUnicodeCharacter{2260}{\ensuremath{\neq}}        % ≠
\DeclareUnicodeCharacter{2264}{\ensuremath{\leq}}        % ≤
\DeclareUnicodeCharacter{2265}{\ensuremath{\geq}}        % ≥
\DeclareUnicodeCharacter{22C5}{\ensuremath{\cdot}}       % ⋅
\DeclareUnicodeCharacter{2295}{\ensuremath{\oplus}}      % ⊕
\DeclareUnicodeCharacter{2297}{\ensuremath{\otimes}}     % ⊗

% ---- Mathematics & symbols ------------------------------------------
\usepackage{amsmath,amssymb,bm,mathrsfs,mathtools}
\usepackage{amsthm}             % theorem environments (load BEFORE hyperref)

% ---- Theorem environments (shared counter set, TECT canonical) -------
% All theorems / lemmas / propositions / corollaries / remarks share a
% single counter so cross-references read consistently across a paper.
% Definitions get a separate counter (definitions are numbered in their
% own sequence by editorial convention).
\newtheorem{theorem}{Theorem}
\newtheorem{lemma}[theorem]{Lemma}
\newtheorem{proposition}[theorem]{Proposition}
\newtheorem{corollary}[theorem]{Corollary}

\theoremstyle{remark}
\newtheorem{remark}[theorem]{Remark}
\newtheorem{example}[theorem]{Example}

\theoremstyle{definition}
\newtheorem{definition}[theorem]{Definition}
\newtheorem{conjecture}[theorem]{Conjecture}
\newtheorem{hypothesis}[theorem]{Hypothesis}

% ---- Tables / floats / graphics --------------------------------------
\usepackage{booktabs}           % \toprule / \midrule / \bottomrule
\usepackage{array}
\usepackage{graphicx}
\usepackage{enumitem}           % \begin{itemize}[leftmargin=...] options
\usepackage{hyphenat}           % \hyp{} for breakable hyphens in \texttt
\usepackage{seqsplit}           % \seqsplit{} for long unbreakable strings
\sloppy                         % allow stretchier inter-word spacing to
                                % reduce overfull \hbox warnings on long URLs
                                % and long \texttt tokens (paragraph-level).
\emergencystretch=3em           % last-resort hbox stretch budget

% ---- Cross-references (hyperref last among the standard packages,
%      cleveref AFTER hyperref) ----------------------------------------
\usepackage[%
  colorlinks=true,%
  linkcolor=blue,%
  citecolor=blue,%
  urlcolor=blue,%
  breaklinks=true,%
  bookmarks=true,%
  bookmarksnumbered=true%
]{hyperref}
\usepackage[capitalise,nameinlink]{cleveref}

% ---- TECT-canonical author block macros -----------------------------
% (defined here so every paper has the same author / affiliation style)
\newcommand{\tectauthor}{%
  \author{Jusang Lee}%
  \email{jtkor@outlook.com}%
  \affiliation{Independent researcher, Republic of Korea}%
}

% ---- TECT-canonical archive footer ----------------------------------
\newcommand{\tectarchivefooter}{%
  \par\medskip
  \noindent\small\textit{Canonical archive}:
  \url{https://github.com/TECT-OpenPhysics/TECT}.\par%
}

% ---- TECT TODO / AUDIT-FLAG / HONEST-SCOPE inline marks --------------
% Used in drafts to highlight pending audit items without disturbing
% the body text style. Suppress via \tectsuppressmarks (define empty).
\providecommand{\tectauditflag}[1]{\textcolor{red}{\textbf{[AUDIT-FLAG: #1]}}}
\providecommand{\tecttodo}[1]{\textcolor{orange}{\textbf{[TODO: #1]}}}
\providecommand{\tecthonestscope}[1]{\textit{Honest scope: #1.}}

% =====================================================================
% End of TECT canonical preamble
% =====================================================================


\begin{document}

\title{Quantum-gates perspective: SO(10) gauge anomalies from the quantum
defect sector}

\author{Jusang Lee}
\email{jtkor@outlook.com}
\affiliation{Independent researcher, Republic of Korea}

\date{\today}

\begin{abstract}
GAP-3 (SO(10) anomaly cancellation) is examined from the
quantum-gates perspective: under the SO(10) embedding (Math157)
derived from the BCC condensate, the spinor representation
$\mathbf{16}$ contains the precise SM fermion content required for
\emph{symbolically exact} cancellation of the five SM perturbative
gauge anomalies and the mixed gravitational
$\mathrm{Tr}(\mathrm{U}(1)_Y)$ anomaly per generation. The
right-handed neutrino singlet $\mathbf{1}(5)$ has SM hypercharge
$Y_{\mathrm{RHN}} = 0$ (Eq.~\eqref{eq:RHN-Y0} below) and therefore
contributes \emph{zero} to any $\mathrm{U}(1)_Y$ trace; the
cancellation is intrinsic to the
$\overline{\mathbf{5}}\oplus\mathbf{10}$ SM matter content alone.
Numerical evaluation cross-checks the symbolic result at machine
precision and is reported as a redundant cross-check, NOT as a
finite-precision estimate. The current canonical tier of GAP-3 is
\textbf{T6 \textsc{proved conditional}} on (H1) the Math157 SO(10)
embedding and on (H2) Pillar~4 hypothesis H1.1 of Math80-AddA;
unconditional promotion is \emph{not} performed in this paper.
\end{abstract}

\maketitle

% =====================================================================
\section{Introduction}
% =====================================================================
The quantum-completion programme (Math60) specifies four closure gates:
GAP-1 (hbar), GAP-2 (BRST), GAP-3 (anomalies), GAP-4 (cosmology). GAP-3
asks: does the quantum defect sector (Pillar 4, gates) support exact
quantum anomaly cancellation?

The answer is yes. The SO(10) spinor representation $\mathbf{16}$ contains
one SM family plus a right-handed neutrino. This precise content is what is
required for anomaly cancellation. No fine-tuning; it is geometrically
forced.

From the quantum-gates perspective, the anomaly cancellation is a
*consequence* of the gauge-group emergence, not an independent constraint.
This paper demonstrates that connection explicitly.

% =====================================================================
\section{Setting and main result}
% =====================================================================
The SO(10) spinor decomposes under
$\mathrm{SU}(5)\times \mathrm{U}(1)_\chi$ as
\begin{equation}
\mathbf{16} = \mathbf{10}(1) \oplus \overline{\mathbf{5}}(-3) \oplus \mathbf{1}(5).
\end{equation}
Further breaking to SM:
\begin{equation}
\mathbf{10}(1) \to (\mathbf{3}, \mathbf{2})_{1/6} \oplus (\overline{\mathbf{3}}, \mathbf{1})_{-2/3},
\quad \overline{\mathbf{5}}(-3) \to (\mathbf{1}, \mathbf{2})_{-1/2} \oplus (\overline{\mathbf{3}}, \mathbf{1})_{1/3}.
\end{equation}

\begin{theorem}[SO(10) anomaly cancellation: symbolic exactness]
\label{thm:main}
Under hypothesis (H1) the SO(10) embedding theorem (Math157), the SO(10)
spinor $\mathbf{16}$ contains the exact SM fermion content required
for symbolic anomaly cancellation. The relevant anomaly coefficients
to verify (under standard SM perturbative-anomaly classification) are:
\begin{itemize}[leftmargin=1.4em]
\item $\mathrm{Tr}(\mathrm{U}(1)_Y^3)$ --- pure hypercharge perturbative.
\item $\mathrm{Tr}(\mathrm{SU}(3)_c^2 \cdot \mathrm{U}(1)_Y)$ --- mixed
  colour--hypercharge.
\item $\mathrm{Tr}(\mathrm{SU}(2)_L^2 \cdot \mathrm{U}(1)_Y)$ --- mixed
  electroweak--hypercharge.
\item $\mathrm{Tr}(\mathrm{U}(1)_Y)$ --- mixed gravitational--hypercharge
  (gravitational anomaly trace).
\end{itemize}
The $\mathrm{SU}(3)_c^3$ anomaly vanishes by vector-like quark
content; the $\mathrm{SU}(2)_L^3$ perturbative anomaly is
group-theoretically trivial, with the relevant constraint being the
Witten global SU(2) anomaly (treated separately). All four
non-trivial coefficients vanish symbolically (exact algebraic zero);
numerical evaluation cross-checks at machine precision. The result
is conditional on Math157 SO(10) embedding; Math281 records the
algebraic-closure observation that the cancellation, once the SO(10)
embedding is fixed, is automatic.
\end{theorem}

% =====================================================================
\section{Proof sketch}
% =====================================================================
\paragraph{Step 1 (SO(10) structure and spinor decomposition) ---}
Math157 derives the complete decomposition of the SO(10) spinor under
all intermediate gauge groups. The decomposition is the standard one,
familiar from GUT literature.

\paragraph{Step 2 (Trace-method computation) ---}
For each anomaly coefficient $\text{Tr}(G^n)$, compute
\begin{equation}
\text{Tr}(G^n) = \sum_f T_f^{(n)} Q_f,
\end{equation}
where the sum is over all fermions in the spinor, and $T_f^{(n)}$, $Q_f$
are Dynkin indices and charges. Math157 tabulates all 36 Dynkin indices
(6 anomaly types × 6 fermion fields) and verifies that the six sums
all equal zero.

\paragraph{Step 3 (Right-handed neutrino role: $Y=0$, NOT $Y=1/2$).}
\textbf{Physics fact (Math314-AddD physics audit, binding).} The
SO(10) right-handed neutrino sits in the $\mathbf{1}(5)$ component
of the
$\mathbf{16} \to \mathbf{10}(1) \oplus \overline{\mathbf{5}}(-3)
\oplus \mathbf{1}(5)$ decomposition under $\mathrm{SU}(5)\times
\mathrm{U}(1)_\chi$. The label ``$(5)$'' attached to the singlet is
the $\mathrm{U}(1)_\chi$ charge in the standard normalisation
$Q_\chi(\mathbf{1}) = +5$; it is \emph{not} the SM hypercharge $Y$,
and it is \emph{not} half the SM hypercharge.

The standard SM hypercharge assignment via
$Y = T_{3R} + \tfrac{1}{2}(B-L)$ on the SO(10)
$\mathbf{16}\supset\mathbf{1}(5)$ singlet gives
\begin{equation}\label{eq:RHN-Y0}
  \boxed{\;
    Y_{\mathrm{RHN}} \;=\; 0
  \;}
\end{equation}
because the right-handed neutrino has $T_{3R} = +\tfrac{1}{2}$ and
$B-L = -1$ (lepton, $L$-conjugate). Therefore the RHN \emph{does not
contribute directly to any SM $\mathrm{U}(1)_Y$ trace}
$\mathrm{Tr}(Y^n)$; it makes \emph{zero} contribution to
$\mathrm{Tr}(\mathrm{U}(1)_Y^3)$,
$\mathrm{Tr}(\mathrm{U}(1)_Y\cdot\mathrm{SU}(2)_L^2)$,
$\mathrm{Tr}(\mathrm{U}(1)_Y\cdot\mathrm{SU}(3)_c^2)$,
and to $\mathrm{Tr}(\mathrm{U}(1)_Y)$. The cancellation of all five
SM perturbative gauge anomalies and of the mixed gravitational
$\mathrm{Tr}(\mathrm{U}(1)_Y)$ anomaly within each generation of the
$\mathbf{16}$ is therefore \emph{intrinsic to the
$\overline{\mathbf{5}}\oplus\mathbf{10}$ SM matter content alone};
the RHN $\mathbf{1}(5)$ does not participate in any of these
cancellations.

The genuine consistency role of the RHN is elsewhere:
(i) algebraic completion of the SO(10) $\mathbf{16}$ multiplet so
that the SO(10) gauge group itself is anomaly-free (SO(10) is
anomaly-free as a simple Lie algebra; the $\mathbf{16}$ is the
spinor that realises this);
(ii) closure of the gauged $\mathrm{U}(1)_{B-L}$ mixed gravitational
anomaly $\mathrm{Tr}(B-L)$ when $B-L$ is gauged in the SO(10)
extension (the RHN carries $(B-L) = +1$ and is required to make
$\mathrm{Tr}(B-L) = 0$ per generation);
(iii) Witten $\mathrm{SU}(2)_L$ global anomaly bookkeeping (the RHN
is an SU(2) singlet, contributes zero doublets, and is therefore
neutral with respect to the Witten constraint).
The Math314-AddD audit identified an earlier draft mis-statement
$Y_{\mathrm{RHN}} = 1/2$, which is incorrect; the corrected
\eqref{eq:RHN-Y0} is the binding canonical value.

\paragraph{Step 4 (Conditional algebraic closure under (H1)+(H2)).}
Math281 records the algebraic-closure observation that, conditional
on (H1) Math157 SO(10) embedding and (H2) Pillar~4 H1.1
(Math80-AddA), the anomaly cancellation is an algebraic identity:
the anomaly coefficients either vanish symbolically (as they do for
the SM $\overline{\mathbf{5}}\oplus\mathbf{10}$ content) or do not
vanish at all; there is no continuous family of solutions. The
verdict therefore inherits the Pillar~4 conditionality and is
recorded at \textbf{T6 \textsc{proved conditional}} on (H1)+(H2).
Promotion to T7 \textsc{proved} (unconditional) requires the
discharge of (H1) and (H2) independently and is \emph{not} performed
in this paper.

% =====================================================================
\section{Symbolic verification (numerically cross-checked)}
% =====================================================================
The relevant anomaly coefficients computed via the spinor-trace
method (Math157) are \emph{symbolically exact zero}; numerical
evaluation cross-checks at machine precision (Math319 V1 acceptance
gate). The coefficient table is:
\begin{align}
\mathrm{Tr}(\mathrm{U}(1)_Y^3) &\equiv 0
  \quad \text{(symbolic exact)} \\
\mathrm{Tr}(\mathrm{SU}(3)_c^2 \cdot \mathrm{U}(1)_Y) &\equiv 0
  \quad \text{(symbolic exact)} \\
\mathrm{Tr}(\mathrm{SU}(2)_L^2 \cdot \mathrm{U}(1)_Y) &\equiv 0
  \quad \text{(symbolic exact)} \\
\mathrm{Tr}(\mathrm{U}(1)_Y) &\equiv 0
  \quad \text{(symbolic exact, gravitational anomaly trace)} \\
\mathrm{Tr}(\mathrm{SU}(3)_c^3) &\equiv 0
  \quad \text{(vector-like quark content)} \\
\mathrm{Tr}(\mathrm{SU}(2)_L^3) &\equiv 0
  \quad \text{(group-theoretically trivial; Witten anomaly}\\
& \quad \text{treated separately as a global anomaly).}
\end{align}
The cancellations are exact at the symbolic level; the numerical
``six-digit precision'' wording of an earlier draft has been
corrected per CLAUDE.md \S15.6 rule \#7 (theorem-vs-estimate notation
discipline) and aligned with Math319 V1 acceptance gate of symbolic
exact zero with numerical cross-check.

% =====================================================================
\section{Discussion --- conditional GAP-3 verdict, current canonical
status of the four GAPs}\label{sec:discussion}
% =====================================================================

From the quantum-gates perspective, Pillar~4 (gauge emergence)
fixes the SO(10) embedding structure. Conditional on (H1) the
Math157 SO(10) embedding theorem (and on Pillar~4 H1.1 of
Math80-AddA), the SM perturbative gauge anomalies and the mixed
gravitational $\mathrm{Tr}(\mathrm{U}(1)_Y)$ anomaly cancel
\emph{symbolically exactly} per generation, by intrinsic
$\overline{\mathbf{5}}\oplus\mathbf{10}$ content of the SO(10)
$\mathbf{16}$, without any direct $Y$-charge contribution from the
RHN singlet (Eq.~\eqref{eq:RHN-Y0}).

\paragraph{Current canonical status of the four GAP gates.}
Per Math307 / Math310 (canonical, post-Math314 audits), the
present canonical tiers are
\[
  \text{GAP-1}=T4, \quad \text{GAP-2}=T6, \quad \text{GAP-3}=T6,
  \quad \text{GAP-4}=T3.
\]
The Stage-2 composite is therefore $\min = T3$ (driven by the
GAP-1 verdict-shell consumption pending in Math312 and the GAP-4
falsifiability-package readiness in Math311). The earlier draft
wording ``GAP-1, GAP-2, and GAP-3 all pass inspection'' is
withdrawn; the correct status is \emph{GAP-2 and GAP-3 conditional
on Pillar~4 H1.1 + companion hypotheses; GAP-1 at T4 with
verdict-shell pending; GAP-4 at T3}.

\paragraph{What the present paper records.}
A symbolically exact, numerically cross-checked verification of the
five SM perturbative gauge anomalies and the
$\mathrm{Tr}(\mathrm{U}(1)_Y)$ gravitational mixed anomaly per
generation, conditional on (H1), with the explicit RHN
hypercharge \eqref{eq:RHN-Y0} fixing the misstatement of an earlier
draft. The paper does \emph{not} certify Stage-2 closure and does
\emph{not} promote any other GAP tier; it is the GAP-3 audit-
auxiliary note recorded against the Math319 V1 verification target.

% =====================================================================
\begin{acknowledgments}
The author thanks Math157 for the SO(10) trace-method formalism, and
Math281 for the algebraic-closure and unconditional-upgrade argument.
This work is part of TECT, canonical archive at
\url{https://github.com/TECT-OpenPhysics/TECT}.
\end{acknowledgments}

% =====================================================================
\bibliographystyle{apsrev4-2}
\bibliography{../_shared/tect-references}

\end{document}
